\section{Relacion de Tareas con Conocimientos Adquiridos Durante los Estudios}

\subsection{Ingenieria del software y arquitectura}
La evolucion de sistemas de estados, descriptores de conflicto y validaciones de integracion permitio aplicar principios de diseno de software trabajados en la titulacion. En particular, la separacion de responsabilidades entre capa de datos, logica de negocio y presentacion, junto con la modularidad para crecimiento progresivo de funcionalidad, se hizo evidente en la reorganizacion de los sistemas de comportamiento. Tambien se reforzo la idea de que la gestion de errores y la observabilidad forman parte del propio diseno del sistema y no solo de la depuracion puntual.

\subsection{Programacion en C\# y desarrollo en Unity}
La mayor parte de la contribucion se implemento en C\# sobre Unity, lo que reforzo competencias practicas relacionadas con el ciclo de vida de componentes y escenas, la integracion entre scripts, prefabs y controladores de animacion, el diagnostico de errores de ejecucion en tiempo real y la construccion y verificacion de builds para plataformas concretas. Este contexto permitio aplicar con criterio los conocimientos de programacion orientada a objetos y de gestion de recursos propios de un motor de juego.

\subsection{Sistemas interactivos y Realidad Virtual}
El trabajo sobre rendimiento, input, tracking y despliegue en Quest 3 conecto directamente con contenidos de sistemas interactivos y tecnicas de programacion grafica. La practica hizo necesaria la identificacion de cuellos de botella, la toma de decisiones de compromiso en entornos VR, la atencion al confort del usuario y la adaptacion de configuraciones de proyecto para hardware objetivo.

\subsection{Redes y comunicacion entre sistemas}
La conexion con servicios web y su robustecimiento permitieron aplicar conceptos de comunicacion entre sistemas. Se trabajo con gestion de mensajes y eventos, validacion de protocolos de intercambio, tratamiento defensivo de entradas no esperadas y priorizacion de trazabilidad para facilitar el mantenimiento evolutivo.

\subsection{Pruebas, calidad y mantenimiento}
La estrategia de aserciones, comprobaciones de estados y verificaciones iterativas reflejo competencias de aseguramiento de calidad. En la practica se abordaron el diseno de casos de prueba orientados a funcionalidad critica, la deteccion temprana de regresiones en cambios de integracion y un enfoque de mejora continua sobre un sistema en produccion interna.

\subsection{Competencias transversales}
Ademas de capacidades tecnicas, la practica reforzo competencias profesionales complementarias. La comunicacion con perfiles de distintas areas, la gestion de prioridades en un backlog tecnico cambiante y la documentacion de decisiones para transferencia de conocimiento al equipo resultaron especialmente relevantes en un proyecto donde conviven programacion, animacion, validacion y configuracion de despliegue.

\subsection{Valor formativo global}
La experiencia de practicas ha servido como puente entre la formacion academica y el desarrollo profesional real en un proyecto VR activo. En terminos formativos, la principal aportacion ha sido aprender a equilibrar tres dimensiones que rara vez aparecen juntas en ejercicios academicos: fiabilidad del sistema, escalabilidad futura y tiempos de entrega del equipo.